\section{Conclusion}

The current iPRA repository already contains the logical kernel, substantial
arithmetic, open-sequent transport, and several mature high-level finite
frameworks.  The most productive next step is therefore not another semantic
layer.  It is a concrete realization of the finite evidence those frameworks
assume.

STRIVE supplies that realization for fixed bounded formulas.  One-child proof
choices are encoded by disjoint prefix sums.  All-child obligations are encoded
by heterogeneous mixed-radix tables.  Positive and negative certificate spaces
are generated together, so failure is an auditable counterstrategy rather than
an opaque Boolean.  The compiler's proofs follow the introduction and
elimination rules of intuitionistic logic, while all numerical loops use
primitive recursion and quantifier-free induction.

The same framework clarifies the embedding of broader mathematics.  A theorem
is presented as a finite observation contract, a demand computation, a
certificate or trace checker, an observer, and an open list of exact residual
producers.  Local certificates are erased, noncanonical witnesses are observed
only through stable outputs, and named choices are introduced only by bounded
least search.  Cyclic finite dependencies are solved by least monotone closure;
compatible runs are solved by backward viability.  Missing moduli, height
bounds, separation estimates, and finite covers remain visible rather than
being attributed to the proof checker.

The result is a framework that is general without being vague.  It has a small
kernel-facing interface, a complete mathematical proof of its central compiler
theorem, executable reference algorithms, and a dependency-ordered Lean route.
Its limitations are equally explicit.  Completing the mixed-radix and bounded-
quantifier modules as checked \(\Prf\;\PRAAx\) derivations would convert a large
class of existing research packets from architectural proposals into reusable
formal infrastructure.
